\subsection{\texorpdfstring{空间 \(L_{F}^{p}\) （ \(1 \leqslant p < +\infty\) ）之间的关系}{空间 L sub F to the p （ 1 <= p < +infinity ）之间的关系}}

对每个实数 \(\alpha > 0\) ，映射 \(z \mapsto |z|^{\alpha-1} \cdot z\) 在 F 中 0 的余集上有定义且连续；此外，由于 \(\left||z|^{\alpha-1} \cdot z\right| = |z|^{\alpha}\) ，这个函数随 z 趋于 0，因而可以连续延拓到点 0，办法是给它在这个点处以值 0，即使 \(\alpha < 1\) 亦然。

定理 7. --- 设 p 与 q 是两个实数，使 \(1 \leqslant p < +\infty\) ，\(1 \leqslant q < +\infty\) 。若函数 f 属于 \(L_{F}^{p}\) ，则函数 \(|f|^{(p/q)-1} \cdot f\) 属于 \(L_{F}^{q}\) ，反之亦然。

按假设，存在具有紧支集的连续函数的序列 \((\mathbf{f}_{n})\) ，使 \(\sum_{n=1}^{\infty}\mathrm{N}_{p}(\mathbf{f}_{n})<+\infty\) 且 \(\mathbf{f}(x)=\sum_{n=1}^{\infty}\mathbf{f}_{n}(x)\) 几乎处处成立（No.~4, 定理 3）。令

\[
\mathbf {g} _ {n} = | \mathbf {f} _ {1} + \mathbf {f} _ {2} + \dots + \mathbf {f} _ {n} | ^ {(p / q) - 1} \cdot (\mathbf {f} _ {1} + \mathbf {f} _ {2} + \dots + \mathbf {f} _ {n});
\]

函数 \(g_{n}\) 具有紧支集且连续；另一方面，

\[
\left| \mathbf {g} _ {n} \right| ^ {q} = \left| \mathbf {f} _ {1} + \mathbf {f} _ {2} + \dots + \mathbf {f} _ {n} \right| ^ {p} \leqslant \left(\sum_ {n = 1} ^ {\infty} \left| \mathbf {f} _ {n} \right|\right) ^ {p} = h ^ {q},
\]

其中数值函数 \(h \geqslant 0\) （有限或非有限）满足不等式

\[
\left(\mathrm{N} _ {q} (h)\right) ^ {q} = \left(\mathrm{N} _ {p} \left(\sum_ {n = 1} ^ {\infty} | \mathbf {f} _ {n} |\right)\right) ^ {p} \leqslant \left(\sum_ {n = 1} ^ {\infty} \mathrm{N} _ {p} (\mathbf {f} _ {n})\right) ^ {p} <   + \infty
\]

此外，\(\mathbf{g}_n(x)\) 几乎处处趋于 \(\mathbf{g}(x) = |\mathbf{f}(x)|^{(p/q)-1} \cdot \mathbf{f}(x)\) ，因此 Lebesgue 定理表明 \(\mathbf{g} \in \mathcal{L}_{\mathrm{F}}^q\) 。逆命题是直接的，因为 \(\mathbf{f} = |\mathbf{g}|^{(q/p)-1} \cdot \mathbf{g}\) 。

可以证明映射 \(\mathbf{f} \mapsto |\mathbf{f}|^{\frac{p}{q}-1} \cdot \mathbf{f}\) 是 \(\mathcal{L}_{\mathrm{F}}^{p}\) 到 \(\mathcal{L}_{\mathrm{F}}^{q}\) 上的一个同胚（§6, 习题 10）。

推论 1. --- 为使函数 \(\mathbf{f}\) 属于 \(\mathcal{L}_{\mathrm{F}}^{p}\) ，必须且只须函数 \(|\mathbf{f}|^{p-1} \cdot \mathbf{f}\) 属于 \(\mathcal{L}_{\mathrm{F}}^{1}\) 。

推论 2. --- 为使正数值函数 f 属于 \(L^{p}\) ，必须且只须 \(f^{p}\) 属于 \(L^{1}\) 。

\pandocbounded{\includegraphics[keepaspectratio]{images/part001_881031c4.jpg}}

注意，若 \(f\) 是任意符号的数值函数，使 \(|f|^p\) 属于 \(\mathcal{L}^1\) ，则 \(f\) 不一定属于 \(\mathcal{L}^p\) （参见 §4, 习题 8）。

